\scanpage{13}

% loc: top continuation
donc il n'y a pas d'isom \uline{canonique} entre $L\otimes L$ et \deleted{\unsure{le}} produit de la famille constante \inserted{sur} $I$ \deleted{\unsure{des|de}} valeur $L$ : il faut, pour choisir un tel isom, choisir un ordre total sur $I$. c'est pire si $I$ est de cardinal $\geq2$, il y a \deleted{\unsure{plus de}} a priori $n!$ isom. possibles entre $\Pi((L_i)_{i\in I})$ et $L^{\otimes 2}$,

% loc: invertible L
Mais si $L$ est inversible, \unsure{la|le} symétrie de $L\otimes L$ \deleted{est} s'identifie à un objet \inserted{$\varepsilon(L)$} du groupe $\operatorname{Aut}(1_C)$, d'ordre 2, et le passage d'un isom. de $\Pi L_i$ avec $L^{\otimes n}$ à un autre, via une permutation $\sigma$ de $[1,n]$, se fait par $\varepsilon(L)^{\operatorname{sign}(\sigma)}$. Il n'y a aucune ambiguité pour le choix dans le seul cas des catégories de Picard [tout objet inv, les $\varepsilon(L)=1$].

% loc: left of inverse discussion
% slanted
\margin{\unsure{kif kif} pour \unsure{pour que l'un des 2|pour un des 2} mais $L\otimes L'\to1$ n'est plus un isom.}

% loc: inverse definition
Inverse d'un élément $L$ inversible : c'est un $L'$ muni d'un isom
\[L\otimes L'\simeq1\]
d'où, en appliquant la symétrie, un \deleted{\unsure{1}} autre isom
\[L'\otimes L\simeq1\]
donc la relation entre $L$ et $L'$ est symétrique, on convient de dire que $(L,L')$ (\deleted{\unsure{sans|non}} ordre fixé) est un couple d'inverses si on se donne un isom
\[\underline{\Pi(L,L')\simeq1}\]

% loc: middle left triangle
\diagram{\begin{tikzcd}[ampersand replacement=\&] L\otimes L'\arrow[d,"s_{L,L'}"']\arrow[dr,"\varphi_{L,L'}"]\&\\L'\otimes L\arrow[r,"\varphi_{L',L}"']\&1\end{tikzcd}}
\uline{\unsure{comm|canon|Comme}}

% loc: middle right
\[L\otimes L'\otimes L\xrightarrow[\sim]{\varphi_{L,L'}\otimes\operatorname{id}_L\simeq\operatorname{id}_L\otimes\psi_{L',L}}L\qquad
\psi_{L',L}:L'\otimes L\simeq1\]
Relation entre $\varphi_{L',L}$ et $\psi_{L',L}$ ? \unsure{Ils|Id.}

% loc: lower formulas
\[
\begin{aligned}
&\varphi_{L,L'}\otimes\operatorname{id}_L\qquad\underbrace{L\otimes L'}\otimes L\longrightarrow L\\
&\operatorname{id}_L\otimes\varphi_{L',L}=\operatorname{id}_L\otimes(\varphi_{L,L'}\circ s_{L',L}):L\otimes L'\otimes L\longrightarrow L
\end{aligned}
\]

% loc: foot left, auxiliary square
$\varphi:M\to1$
\diagram{\begin{tikzcd}[ampersand replacement=\&,column sep=small] L\otimes M\arrow[rr,"s_{L,M}"]\arrow[d,"L\otimes\varphi"']\&\&M\otimes L\arrow[d,"\varphi\otimes L"]\\L\otimes1\arrow[rr,"{\unsure{s|k}_{L,1}}"]\arrow[dr,"\simeq"']\&\&1\otimes L\arrow[dl,"\simeq"]\\\&L\&\end{tikzcd}}

% loc: foot centre, two contractions
\diagram{\begin{tikzcd}[ampersand replacement=\&,column sep=tiny,row sep=small]
\underbrace{L\qquad L'}\arrow[d,"\varphi"]\& L\&\quad L\&\underbrace{L'\qquad L}\arrow[d]\\
1\&\& L\&\underbrace{L^{\unsure{*|a}}\qquad L'^{\unsure{*|a}}}\arrow[d,"\varphi"]\\
\&\&\&1
\end{tikzcd}}

% loc: foot right, bracket and array
\diagram{$\left]\ \begin{array}{cccc}&L&L'&L\\ &L&L^{\unsure{\xi|*}}&L'\\ \unsure{\xi|\varepsilon}&L&L'&L\end{array}\right.$}

% archive sheet number: 12
